% !PRINTED 358
\weberpage{358}
%<<CONT
linear function containing three arbitrary constants, it
follows from (1):
\[
f = x_1 (V + 2pu + x_1p^2) - (u + px_1)^2,
\]
which is again of the form (1).

If now $y_1 = 0$ is the equation of a second double tangent distinct from $x_1$, then we can choose the constants in $p$ (and indeed in infinitely many different ways) so that the conic section $u + px_1 = 0$ passes through the two points of contact of $y_1$, and we may therefore assume that the function $u$ has already been so chosen in the form (1). Then the cubic form $V$ must also vanish at the same points.

But still more: since the line $y_1 = 0$ is to be a double tangent, it must be the case that, if we introduce $x_1, y_1$ and any third linear function $z$ independent of them as variables, at the points of contact
\[
f'(x_1) = 0, \quad f'(z) = 0
\]
and from this it follows by (1):
\[
V'(x_1) = 0, \quad V'(z) = 0,
\]
that is, the curve of the third order $V = 0$ is touched by the line $y_1 = 0$ at the points of contact with $f$, and this is possible only if $V$ splits and contains the function $y_1$ as a divisor. Accordingly we obtain a new shape of the function $f$:
\weq{2}{f = x_1 y_1 v - u^2,}
in which $u, v$ are functions of the second degree.

If $x_1, y_1$ are any arbitrary linear forms, and $u, v$ quadratic forms of three variables, then the function $f$ determined by (2), set equal to zero, represents a curve of the fourth order, of which $x_1$ and $y_1$ are two double tangents.

If conversely $x_1$ and $y_1$ are two arbitrary double tangents of a curve of the fourth order $f = 0$, then $f$ can be brought to the form (2).

Let us imagine the coefficients of $x_1$ and $y_1$ as variable, and let these quantities vary in such a way that $x_1$ and $y_1$ intersect on the curve $f$; then a double point will occur.

The double tangents degenerate into two tangents issuing from the double point.
